\subsection*{Part (a). Initialization against training length}

\lead{Setup.} Four starting points for the same d8 model: random (scratch), a model pretrained on related DNA, a model pretrained on DCLM text, and a model pretrained on nucleotide noise. Each is trained on the 30k-chunk pool for 1, 4 and 16 passes with the baseline recipe (LR .003, WD 3.2, warmup + linear decay to zero at the end of the run), plus lower LRs (.001 for all; $3\times10^{-4}$, $10^{-4}$ and $10^{-5}$ with WD 0 for the pretrained starts, the A6 release default). Table~\ref{tab:p4a} reports the best LR per cell.

\lead{Findings} (Table~\ref{tab:p4a}).

\begin{itemize}
\item Pretraining on related DNA is worth most when training is short and gentle: one pass at LR $3\times10^{-4}$ (WD 0, the release default) reaches 1.2281, the best single model in this report and 0.035 below the supplied 32-pass baseline; at the scratch LR .003 the same pass gives 1.269, against 1.299 from scratch, 1.294 from DCLM and 1.312 from noise. Four gentle passes still give 1.2412.
\item The gap closes fast. After four passes scratch is at 1.276 and related at 1.270; DCLM (1.275) is indistinguishable from scratch and noise (1.283) is still behind.
\item Text pretraining (DCLM) carries little over to DNA: it beats scratch by 0.005 at one pass and by nothing at four. Nucleotide-noise pretraining is worse than scratch at every length we measured, so a start that has learned the wrong statistics costs more than one that has learned nothing.
\item At 16 passes the order inverts: scratch (1.2638) and DCLM (1.2637) now edge out the related start (1.2653), and noise (1.2676) is still last. Full fine-tuning at the scratch LR (.003) overwrites the pretrained features by the time the data have been seen 16 times; the related start never improves much beyond its 1-pass value, at LR .003 or .001 (Part (c)); from scratch, LR .001 is the better 16-pass choice (1.2618 against 1.2638). At LR $3\times10^{-4}$ with WD 0 the related start collapses to 1.64 by pass 16 (1.2295 at pass 1), so the gentle recipe needs a short horizon or a regularizer; the strong initialization is an advantage only while training is short, gentle, or constrained (Part (c)).
\end{itemize}
